\section{Baselines}
\label{sec:dataset-baselines}

\subsection{Text-Only Classifiers}
\label{subsec:text-only-baselines}

Section~\ref{subsec:lexical-separability} requires the dataset to be difficult for classifiers that take only the request text as input, and it names three text representations that these classifiers must be trained on:

\begin{itemize}
    \item \textbf{bag-of-words}: a logistic regression over raw word counts (unigrams and bigrams);
    \item \textbf{TF-IDF}: a logistic regression over the same counts, weighted by TF-IDF;
    \item \textbf{sentence embeddings}: a logistic regression over sentence embeddings.
\end{itemize}

For each classifier, \emph{accuracy} is measured, given by
\begin{equation}
    \text{accuracy} = \frac{\text{requests classified correctly}}{\text{all requests}},
\end{equation}
where a request is classified correctly when the classifier's predicted sub-agent matches the correct sub-agent for that request.

Accuracy is measured with \emph{5-fold cross-validation}: the dataset is split into five folds, each classifier is trained on four folds and evaluated on the remaining fold, and this is repeated five times so that every fold serves once as the validation fold. Each fold, when it acts as the validation fold, produces a predicted sub-agent for every request in it. This is the \emph{out-of-fold prediction} for that request, since the request was never in the training data used to produce it. Collecting the out-of-fold prediction from each of the five folds gives one predicted sub-agent for every request in the dataset. Accuracy is then computed once, as defined above, by comparing each request's out-of-fold prediction against its correct sub-agent. This procedure uses every sample for both training and validation, without letting the same sample serve as both at once, and it prevents the reported accuracy from depending on one arbitrary split of the data. Folds are grouped by scenario, so a scenario's samples never split between the training fold and the validation fold. This matters for an \texttt{adversarial} pair, since its two requests share nearly all their wording: if one request were used for training and the other for validation, the classifier could match the validation request to its near-duplicate seen during training, and score high without learning to separate the sub-agents.

The sentence embeddings are obtained with \texttt{BAAI/bge-small-en-v1.5}, an embedding model that maps a request to a fixed-size vector encoding its meaning, so that requests with a similar meaning are mapped to nearby vectors.

Each classifier is trained and evaluated four times, on four different subsets of the samples:

\begin{itemize}
    \item \textbf{All}: the whole dataset;
    \item \textbf{Plain}: only the \texttt{plain} tier;
    \item \textbf{Natural}: only the \texttt{natural} tier;
    \item \textbf{Adversarial}: only the \texttt{adversarial} tier.
\end{itemize}

The accuracy obtained on both the single-turn and the multi-turn dataset is reported in Table~\ref{tab:text-baselines}.

\begin{table}[h]
    \centering
    \begin{tabular}{@{}l|lll|lll@{}}
        \toprule
        & \multicolumn{3}{c|}{\textbf{Single-turn}} & \multicolumn{3}{c}{\textbf{Multi-turn}} \\
        \textbf{Tier} & \textbf{Bag-of-words} & \textbf{TF-IDF} & \textbf{Embedding} & \textbf{Bag-of-words} & \textbf{TF-IDF} & \textbf{Embedding} \\
        \midrule
        Plain       & 0.900 & 0.993 & 0.927 & 0.840 & 0.987 & 0.873 \\
        Natural     & 0.708 & 0.730 & 0.653 & 0.665 & 0.677 & 0.592 \\
        Adversarial & 0.631 & 0.566 & 0.478 & 0.599 & 0.508 & 0.411 \\
        \midrule
        All         & 0.684 & 0.664 & 0.582 & 0.645 & 0.613 & 0.519 \\
        \bottomrule
    \end{tabular}
    \caption[Text-only classifier accuracy by difficulty tier, both datasets]{Accuracy of the three text-only classifiers, per difficulty tier, on the single-turn and the multi-turn dataset.}
    \label{tab:text-baselines}
\end{table}

On the whole dataset, no classifier reaches an accuracy above 0.69, on either the single-turn or the multi-turn dataset. Broken down by tier, accuracy falls from \texttt{plain} to \texttt{natural} and falls again from \texttt{natural} to \texttt{adversarial}, for all three classifiers, on both datasets. This is the decreasing accuracy that the tiers are designed to produce (Section~\ref{subsec:difficulty-tiers}). On \texttt{plain}, the request states only the binding constraint, so its vocabulary alone gives away the correct sub-agent. On \texttt{natural}, the request also states a waived constraint, so the vocabulary of a second sub-agent appears in the same request. On \texttt{adversarial}, the request is one half of a pair that shares nearly all its wording with its twin, so the words alone cannot tell which of the two constraints is binding.

The intrinsic difficulty of the \texttt{adversarial} design is also confirmed by the TF-IDF classifier results, which are outperformed by a non-TF-IDF-weighted bag-of-words classifier only on that tier, on both datasets. The two requests in an adversarial pair share nearly all their wording, and this high overlap makes TF-IDF down-weight words that still carry part of the routing signal.

The embedding classifier beats bag-of-words only on the \texttt{plain} tier. On \texttt{plain}, the request states one constraint, and the embedding captures it directly. On \texttt{natural}, the request also states a second, waived constraint. On \texttt{adversarial}, the request is nearly identical to its pair. In both cases, the embedding represents the whole request, and most of the request is the shared scenario text rather than the constraint that decides the sub-agent. Two requests built from the same scenario can then end up with similar embeddings even when they belong to different sub-agents, which lowers accuracy.

\subsection{Vocabulary and Routability}
\label{subsec:baseline-vocabulary}

Table~\ref{tab:baseline-terms} shows the single-turn TF-IDF terms with the highest weight, for each sub-agent. This confirms that most of the accuracy on \texttt{natural} and \texttt{adversarial} comes from requirement-specific vocabulary.

\begin{table}[h]
    \centering
    \begin{tabular}{@{}ll@{}}
        \toprule
        \textbf{Sub-agent} & \textbf{Highest-weighted TF-IDF terms} \\
        \midrule
        \texttt{legacy\_compat}     & contract, API, clients, existing, without breaking, published \\
        \texttt{time\_to\_market}       & weeks, working, demo, two weeks, ship, product \\
        \texttt{resource\_opt}      & memory, budget, ceiling, CPU, cloud, under \\
        \texttt{scale\_arch}        & scale, infrastructure, shared, horizontally, decouple, across \\
        \texttt{strict\_compliance} & audit, data, compliance, logging, security, officer \\
        \bottomrule
    \end{tabular}
    \caption[Highest-weighted TF-IDF terms per sub-agent]{Highest-weighted TF-IDF terms per sub-agent.}
    \label{tab:baseline-terms}
\end{table}

Section~\ref{subsec:lexical-separability} requires every request to stay routable to one sub-agent and to read as natural. The text-only baseline accuracies in Table~\ref{tab:text-baselines} are the lowest achievable under that requirement. To lower this accuracy further, the highest-weighted terms in Table~\ref{tab:baseline-terms} could be banned from generation, so that the generator cannot rely on them. The result is that the generator relies on synonyms instead, and a synonym is equally separable for a classifier. Banning every synonym as well introduces ambiguity: requests stop reading as natural and become hard to route even for a human.

\subsection{Difficulty Beyond the Baselines}
\label{subsec:encoder-robustness}

The generation pipeline is designed to minimize the accuracy of the three baselines whose results are reported in Table~\ref{tab:text-baselines}. Their low accuracy therefore does not show that the datasets are also difficult for other text-only classifiers. To show this, two more logistic regressions are trained on the representations of two different encoders, on the same scenario-grouped folds (Table~\ref{tab:modernbert-baseline}).

\begin{table}[H]
    \centering
    \small
    \begin{tabular}{@{}l|cc|cc@{}}
        \toprule
        & \multicolumn{2}{c|}{\textbf{Single-turn}} & \multicolumn{2}{c}{\textbf{Multi-turn}} \\
        \textbf{Tier} & \textbf{ModernBERT} & \textbf{BGE-M3} & \textbf{ModernBERT} & \textbf{BGE-M3} \\
        \midrule
        Plain       & 0.827 & 0.967 & 0.747 & 0.940 \\
        Natural     & 0.665 & 0.770 & 0.607 & 0.643 \\
        Adversarial & 0.552 & 0.524 & 0.462 & 0.431 \\
        \midrule
        All         & 0.618 & 0.654 & 0.541 & 0.555 \\
        \bottomrule
    \end{tabular}
    \caption[ModernBERT-large and BGE-M3 accuracy by difficulty tier]{Accuracy of the logistic regressions over ModernBERT-large representations and over BGE-M3 sentence embeddings, per difficulty tier, on the single-turn and the multi-turn dataset.}
    \label{tab:modernbert-baseline}
\end{table}

The first is trained on the representations of ModernBERT-large \cite{warnerSmarterBetterFaster2024}. ModernBERT-large maps each token of the request to a vector, and the representation of the request is obtained by mean pooling these vectors. The second is trained on the sentence embeddings of \texttt{BAAI/bge-m3}, a larger embedding model than \texttt{BAAI/bge-small-en-v1.5}.

Neither classifier reaches an accuracy above 0.69 on the whole dataset, single-turn or multi-turn, and the accuracy of both falls from \texttt{plain} to \texttt{natural} to \texttt{adversarial}, as for the three baselines. The difficulty is therefore a property of the datasets and is not specific to the three classifiers chosen as baselines.
